\documentclass[pdf,blends,slideColor,colorBG]{prosper}

\usepackage{moreverb}
\usepackage{macros}

\begin{document}

\title{Le mod�le Entit�/Association}
\author{Philippe Rigaux}
\email{rigaux@lri.fr}
\institution{LRI}

%------------------------------------------------------------
\part{Conception d'une \\base de donn�es\\relationnelle}
%------------------------------------------------------------


\begin{tsp}{Quelques principes}{

On cherche � d�finir une \guill{bonne} organisation
des donn�es. Quelques crit�res\,:
\begin{itemize}
   \item Pas de redondance (une information figure une et une seule fois)

  \item Pas de perte  (on est toujours capable
          de retrouver une information)

   \item On repr�sente ce qui nous int�resse, ni plus, ni moins.
\end{itemize}

Pas si facile en pratique.
}\end{tsp}
%%%%%%%%%%%%%%%%%%%%%%%%%%%%%%%%%%%%%%%%%%%%%%%%%%%
\begin{tsp}{Illustration}{

\begin{center}
\begin{small}
\begin{tabular}{l|c|l|l|c}
 titre &  ann�e &  nomMES &  pr�nomMES &  ann�eNaiss \\
\hline
Alien & 1979 & Scott & Ridley & 1943\\
Vertigo  & 1958 & Hitchcock & Alfred & 1899\\
Psychose & 1960 & Hitchcock & Alfred & 1899\\
Kagemusha & 1980 & Kurosawa & Akira & 1910\\
Volte-face & 1997 & Woo & John & 1946\\
Pulp Fiction & 1995 & Tarantino & Quentin & 1963 \\
Titanic & 1997 & Cameron & James & 1954\\
Sacrifice & 1986 & Tarkovski & Andrei &1932\\
\hline
\end{tabular}
\end{small}
\end{center}

}\end{tsp}
%%%%%%%%%%%%%%%%%%%%%%%%%%%%%%%%%%%%%%%%
\begin{tsp}{Pourquoi �a ne va pas}{

Une seule table, �a semble pratique mais �a soul�ve beaucoup
de probl�mes.

\begin{itemize}
   \item Redondance (plusieurs fois le m�me film\,?!)

   \item Perte d'information\,: je supprime Volte-face,
     et John Woo dispara�t par la m�me occasion.

   \item Toutes sortes d'anomalies et d'incoh�rences peuvent appara�tre.
\end{itemize}

Question\,: qu'est-ce qu'un film, qu'un metteur en sc�ne, comment
sont-ils li�s, etc.
}\end{tsp}
%%%%%%%%%%%%%%%%%%%%%%%%%%%%%%%%%%%%%%%%
\begin{tsp}{Repr�sentation ind�pendante}{

\noindent
\begin{small}
\begin{center}
\begin{minipage}[b]{4cm}
\begin{center}
\begin{small}
\begin{tabular}{l|c}
{\bf titre} & ann�e  \\
\hline
Alien & 1979 \\
Vertigo  & 1958 \\
Psychose & 1960 \\
Kagemusha & 1980 \\
Volte-face & 1997 \\
Pulp Fiction & 1995  \\
Titanic & 1997 \\
Sacrifice & 1986 \\
\hline
\end{tabular}
\end{small}

La table des films
\end{center}
\end{minipage} \
\begin{minipage}[b]{5cm}
\begin{center}
\begin{tabular}{l|l|l|c}
{\bf id} &  nomMES & pr�nomMES &  ann�eNaiss \\
\hline
1 & Scott & Ridley & 1943\\
2 & Hitchcock & Alfred & 1899\\
3 & Kurosawa & Akira & 1910\\
4 & Woo & John & 1946\\
5 & Tarantino & Quentin & 1963 \\
6 & Cameron & James & 1954\\
7 & Tarkovski & Andrei &1932\\
\hline
\end{tabular}

La table des r�alisateurs
\end{center}
\end{minipage}
\end{center}
\end{small}

}\end{tsp}
%%%%%%%%%%%%%%%%%%%%%%%%%%%%%%%%%%%%%%%%
\begin{tsp}{Repr�sentation ind�pendante \textbf{et} li�e}{
\begin{small}
\begin{center}
\begin{minipage}[b]{3cm}
\begin{tabular}{l|c|c}
\cle{titre} & ann�e  & \fk{idMES} \\
\hline
Alien & 1979 & 1 \\
Vertigo  & 1958 & 2\\
Psychose & 1960 & 2 \\
Kagemusha & 1980 & 3\\
Volte-face & 1997 & 4 \\
Pulp Fiction & 1995 &  5 \\
Titanic & 1997 & 6 \\
Sacrifice & 1986 & 7\\
\hline
\end{tabular}

\centerline{La table des films}
\end{minipage} \hspace*{1cm} \
\begin{minipage}[b]{5cm}
\begin{center}
\begin{tabular}{l|l|l|c}
{\bf id} &  nomMES & pr�nomMES &  ann�eNaiss \\
\hline
1 & Scott & Ridley & 1943\\
2 & Hitchcock & Alfred & 1899\\
3 & Kurosawa & Akira & 1910\\
4 & Woo & John & 1946\\
5 & Tarantino & Quentin & 1963 \\
6 & Cameron & James & 1954\\
7 & Tarkovski & Andrei &1932\\
\hline
\end{tabular}

La table des r�alisateurs
\end{center}
\end{minipage}
\end{center}
\end{small}

}\end{tsp}

%%%%%%%%%%%%%%%%%%%%%%%%%%%%%%%%%%%%%%%%
\begin{tsp}{Vue d'ensemble}{

\figSimple{films}

}\end{tsp}
%%%%%%%%%%%%%%%%%%%%%%%%%%%%%%%%%%%%%%%%
\begin{tsp}{Concepts\,: entit�s et type d'entit�s}{

Une entit�\,:

\begin{itemize}
   \item Une identit�

   \item des attributs

   \item appartient � un {\em groupe} 
\end{itemize}

Entit� = instance d'un type d'entit�
}\end{tsp}
%%%%%%%%%%%%%%%%%%%%%%%%%%%%%%%%%%%%%%%%
\begin{tsp}{Types d'entit�s}{

\figSimple{entite}

}\end{tsp}
%%%%%%%%%%%%%%%%%%%%%%%%%%%%%%%%%%%%%%%%
\begin{tsp}{Graphe d'une association}{

\figSimple{graphe-ea}

}\end{tsp}
%%%%%%%%%%%%%%%%%%%%%%%%%%%%%%%%%%%%%%%%
\begin{tsp}{Association}{

\figSimple{assoc}

}\end{tsp}
%%%%%%%%%%%%%%%%%%%%%%%%%%%%%%%%%%%%%%%%
\begin{tsp}{Autre graphe}{

\figSimple{graphe-ea2}

}\end{tsp}
%%%%%%%%%%%%%%%%%%%%%%%%%%%%%%%%%%%%%%%%
\begin{tsp}{Autre association}{

\figSimple{assoc2}

}\end{tsp}

\end{document}
